\documentclass[tikz,border=3mm,multi=tikzpicture]{standalone}
\usepackage{eleve-fisica}

\begin{document}

% FIGURA 1 — fig-01-troca-entre-forca-e-distancia
\begin{tikzpicture}[fisica figura, x=1cm, y=1cm]
  \path[use as bounding box] (-0.2,-0.2) rectangle (10.7,9.3);
  \node[font=\Large\bfseries, text=eleveazul] at (5.15,8.75)
    {mesma carga, mesma altura};
  \draw[fisica superficie] (0.8,1.05) -- (9.65,1.05);
  \draw[fisica superficie] (8.8,1.05) -- (8.8,6.4);
  \node[fisica corpo, minimum width=1.45cm, minimum height=1.0cm] (A)
    at (2.0,5.85) {};
  \draw[fisica forca] (A.center) -- ++(0,2.0)
    node[fisica rotulo, above] {força maior};
  \draw[decorate, decoration={brace, amplitude=7pt}, elevecinza, line width=1pt]
    (1.0,1.05) -- (1.0,5.85)
    node[midway, left=12pt, fisica rotulo] {caminho curto};

  \fill[elevecinzaclaro] (4.0,1.05) -- (8.8,1.05) -- (8.8,6.4) -- cycle;
  \draw[fisica superficie] (4.0,1.05) -- (8.8,6.4);
  \begin{scope}[shift={(6.55,3.9)}, rotate=48]
    \node[fisica corpo, minimum width=1.45cm, minimum height=1.0cm] (B) at (0,0) {};
    \draw[fisica forca] (B.center) -- ++(2.2,0)
      node[fisica rotulo, right] {força menor};
  \end{scope}
  \node[fisica rotulo, rotate=48] at (5.45,3.25) {caminho longo};
\end{tikzpicture}

% FIGURA 2 — fig-02-elementos-e-classes-de-alavancas
\begin{tikzpicture}[fisica figura, x=1cm, y=1cm]
  \path[use as bounding box] (-0.2,-1.1) rectangle (10.6,12.4);
  \node[font=\Large\bfseries, text=eleveazul] at (5.1,11.85)
    {interfixa: apoio no meio};
  \draw[fisica superficie] (1.0,9.45) -- (9.2,9.45);
  \fill[elevecinza] (5.1,9.45) -- (4.55,8.55) -- (5.65,8.55) -- cycle;
  \draw[fisica forca] (2.0,9.45) -- ++(0,-1.35)
    node[fisica rotulo, below] {$\vec F$};
  \draw[fisica forca] (8.2,9.45) -- ++(0,-1.35)
    node[fisica rotulo, below] {$\vec R$};
  \node[fisica medida] at (5.1,8.15) {$A$};

  \node[font=\Large\bfseries, text=eleveazul] at (5.1,7.1)
    {inter-resistente: carga no meio};
  \draw[fisica superficie] (1.0,4.75) -- (9.2,4.75);
  \fill[elevecinza] (2.0,4.75) -- (1.45,3.85) -- (2.55,3.85) -- cycle;
  \draw[fisica forca] (5.1,4.75) -- ++(0,-1.35)
    node[fisica rotulo, below] {$\vec R$};
  \draw[fisica forca] (8.4,4.75) -- ++(0,1.35)
    node[fisica rotulo, above] {$\vec F$};
  \node[fisica medida] at (2.0,3.45) {$A$};

  \node[font=\Large\bfseries, text=eleveazul] at (5.1,2.75)
    {interpotente: força no meio};
  \draw[fisica superficie] (1.0,0.5) -- (9.2,0.5);
  \fill[elevecinza] (2.0,0.5) -- (1.45,-0.4) -- (2.55,-0.4) -- cycle;
  \draw[fisica forca] (5.1,0.5) -- ++(0,1.35)
    node[fisica rotulo, above] {$\vec F$};
  \draw[fisica forca] (8.4,0.5) -- ++(0,-1.0)
    node[fisica rotulo, below] {$\vec R$};
  \node[fisica medida] at (2.0,-0.8) {$A$};
\end{tikzpicture}

% FIGURA 3 — fig-03-polia-fixa-e-polia-movel
\begin{tikzpicture}[fisica figura, x=1cm, y=1cm]
  \path[use as bounding box] (-0.2,-0.3) rectangle (11.0,10.2);
  \node[font=\Large\bfseries, text=eleveazul] at (2.75,9.65) {polia fixa};
  \draw[fisica superficie] (0.65,8.7) -- (4.9,8.7);
  \draw[elevecinza, line width=1.4pt] (2.75,8.7) -- (2.75,7.85);
  \draw[elevecinza, fill=white, line width=1.4pt] (2.75,7.2) circle (0.65);
  \fill[elevecinza] (2.75,7.2) circle (2.5pt);
  \draw[elevecinza, line width=1.5pt] (2.1,7.2) -- (2.1,3.7);
  \draw[elevecinza, line width=1.5pt] (3.4,7.2) -- (3.4,3.7);
  \node[fisica corpo, minimum width=1.5cm, minimum height=1.05cm] at (2.1,3.15) {carga};
  \draw[fisica forca] (3.4,3.7) -- ++(0,-1.7)
    node[fisica rotulo, below] {$\vec F$};
  \node[fisica rotulo, align=center] at (2.75,1.05) {muda o\\sentido};

  \node[font=\Large\bfseries, text=eleveazul] at (8.05,9.65) {polia móvel};
  \draw[fisica superficie] (5.85,8.7) -- (10.25,8.7);
  \fill[elevecinza] (6.7,8.7) circle (3pt);
  \draw[elevecinza, line width=1.5pt] (6.7,8.7) -- (6.7,5.75)
    arc[start angle=180,end angle=0,radius=1.05] -- (8.8,8.1);
  \draw[elevecinza, fill=white, line width=1.4pt] (7.75,5.75) circle (1.05);
  \fill[elevecinza] (7.75,5.75) circle (2.5pt);
  \draw[elevecinza, line width=1.4pt] (7.75,4.7) -- (7.75,3.75);
  \node[fisica corpo, minimum width=1.5cm, minimum height=1.05cm] at (7.75,3.2) {carga};
  \draw[fisica forca] (8.8,8.1) -- ++(0,-1.65)
    node[fisica rotulo, right] {$\vec F$};
  \node[fisica medida] at (6.25,5.8) {$T$};
  \node[fisica medida] at (9.25,5.8) {$T$};
  \node[fisica rotulo, align=center] at (7.75,1.05) {dois trechos\\sustentam};
\end{tikzpicture}

% FIGURA 4 — fig-04-trechos-que-sustentam-a-carga
\begin{tikzpicture}[fisica figura, x=1cm, y=1cm]
  \path[use as bounding box] (-0.2,-0.2) rectangle (10.8,9.9);
  \draw[fisica superficie] (0.8,8.85) -- (9.85,8.85);
  \foreach \x in {2.0,3.9,5.8,7.7}
    \draw[elevecinza, line width=1.5pt] (\x,8.85) -- (\x,4.55);
  \foreach \x in {2.95,6.75}
    \draw[elevecinza, fill=white, line width=1.4pt] (\x,4.55) circle (0.95);
  \draw[elevecinza, line width=1.5pt]
    (2.0,4.55) arc[start angle=180,end angle=0,radius=0.95];
  \draw[elevecinza, line width=1.5pt]
    (5.8,4.55) arc[start angle=180,end angle=0,radius=0.95];
  \draw[elevecinza, line width=1.4pt] (2.95,3.6) -- (2.95,2.8);
  \draw[elevecinza, line width=1.4pt] (6.75,3.6) -- (6.75,2.8);
  \draw[fisica superficie] (2.95,2.8) -- (6.75,2.8);
  \node[fisica corpo, minimum width=3.0cm, minimum height=1.35cm] at (4.85,1.75) {carga};
  \foreach \x in {2.0,3.9,5.8,7.7}
    \node[fisica medida] at (\x,6.75) {$T$};
  \node[font=\Large\bfseries, text=eleveazul] at (4.85,9.45)
    {quatro trechos sustentam};
  \draw[fisica forca] (9.15,8.2) -- ++(0,-2.1)
    node[fisica rotulo, below] {$\vec F$};
  \node[fisica rotulo] at (4.85,0.35) {$4T=P$ no modelo ideal};
\end{tikzpicture}

% FIGURA 5 — fig-05-rampas-com-mesma-altura
\begin{tikzpicture}[fisica figura, x=1cm, y=1cm]
  \path[use as bounding box] (-0.2,-0.2) rectangle (10.8,9.4);
  \node[font=\Large\bfseries, text=eleveazul] at (5.15,8.8)
    {mesma altura};
  \draw[fisica superficie] (0.75,0.75) -- (9.75,0.75);
  \draw[fisica superficie] (9.0,0.75) -- (9.0,6.9);
  \draw[fisica superficie] (2.0,0.75) -- (9.0,6.9);
  \draw[fisica superficie] (5.35,0.75) -- (9.0,6.9);
  \node[fisica rotulo, rotate=41] at (5.0,4.25) {caminho maior};
  \node[fisica rotulo, rotate=59] at (7.45,4.35) {caminho menor};
  \draw[fisica forca] (4.15,2.65) -- ++(1.35,1.18);
  \node[fisica rotulo] at (2.9,5.05) {força menor};
  \draw[fisica forca] (6.95,3.5) -- ++(0.85,1.45);
  \node[fisica rotulo, anchor=west] at (8.2,4.95) {força maior};
  \draw[elevelaranja, line width=1.1pt] (2.9,0.75) arc[start angle=0,end angle=41,radius=0.9];
  \draw[elevelaranja, line width=1.1pt] (6.1,0.75) arc[start angle=0,end angle=59,radius=0.75];
\end{tikzpicture}

% FIGURA 6 — fig-06-roda-cunha-e-parafuso
\begin{tikzpicture}[fisica figura, x=1cm, y=1cm]
  \path[use as bounding box] (-0.2,-0.2) rectangle (10.8,12.0);
  \node[font=\Large\bfseries, text=eleveazul] at (5.15,11.45) {roda e eixo};
  \draw[eleveazul, fill=eleveazulclaro, line width=1.4pt] (5.15,9.2) circle (1.35);
  \draw[elevecinza, fill=white, line width=1.3pt] (5.15,9.2) circle (0.48);
  \draw[fisica movimento] (6.35,9.85) arc[start angle=30,end angle=-65,radius=1.4];
  \node[fisica medida] at (7.55,9.2) {giro};

  \node[font=\Large\bfseries, text=eleveazul] at (5.15,7.15) {cunha};
  \fill[elevelaranjaclaro] (3.4,4.7) -- (6.9,4.7) -- (5.15,6.55) -- cycle;
  \draw[fisica superficie] (3.4,4.7) -- (5.15,6.55) -- (6.9,4.7);
  \draw[fisica forca] (5.15,6.5) -- ++(0,-1.25)
    node[fisica rotulo, right] {$\vec F$};
  \draw[fisica movimento] (4.55,5.15) -- ++(-1.25,0);
  \draw[fisica movimento] (5.75,5.15) -- ++(1.25,0);

  \node[font=\Large\bfseries, text=eleveazul] at (5.15,3.65) {parafuso};
  \draw[elevecinza, line width=1.5pt] (4.35,0.45) -- (4.35,3.05);
  \draw[elevecinza, line width=1.5pt] (5.95,0.45) -- (5.95,3.05);
  \foreach \y in {0.7,1.25,1.8,2.35,2.9}
    \draw[eleveazul, line width=1.4pt] (4.05,\y) -- (6.25,\y+0.35);
  \draw[fisica movimento] (6.55,2.45) arc[start angle=45,end angle=-110,radius=0.85];
  \draw[fisica movimento] (7.45,1.5) -- ++(0,-1.0)
    node[fisica rotulo, below] {avanço};
\end{tikzpicture}

\end{document}
